\documentclass[aspectratio=1610,slidestop]{beamer}

\author{Mihai Lefter}
\title{Python Programming}
\providecommand{\mySubTitle}{Data Types}
\providecommand{\myConference}{Programming Course}
\providecommand{\myDate}{27-11-2018}
\providecommand{\myGroup}{}
\providecommand{\myDepartment}{}
\providecommand{\myCenter}{}

\usetheme{lumc}

\usepackage{minted}
\usepackage{tikz}
\usepackage[many]{tcolorbox}

\definecolor{monokaibg}{HTML}{272822}
\definecolor{emailc}{HTML}{1e90FF}


\newenvironment{ipython}
 {\begin{tcolorbox}[title=IPython,
                   title filled=false,
                   fonttitle=\scriptsize,
                   fontupper=\footnotesize,
                   enhanced,
                   colback=monokaibg,
                   drop small lifted shadow,
                   boxrule=0.1mm,
                   left=0.1cm,
                   arc=0mm,
                   colframe=black]}
 {\end{tcolorbox}}


\newenvironment{terminal}
 {\begin{tcolorbox}[title=terminal,
                   title filled=false,
                   fonttitle=\scriptsize,
                   fontupper=\footnotesize,
                   enhanced,
                   colback=monokaibg,
                   drop small lifted shadow,
                   boxrule=0.1mm,
                   left=0.1cm,
                   arc=0mm,
                   colframe=black]}
 {\end{tcolorbox}}


\newcommand{\hrefcc}[2]{\textcolor{#1}{\href{#2}{#2}}}
\newcommand{\hrefc}[3]{\textcolor{#1}{\href{#2}{#3}}}

\newcounter{cntr}
\renewcommand{\thecntr}{\texttt{[\arabic{cntr}]}}

\newenvironment{pythonin}[1]
{\VerbatimEnvironment
  \begin{minipage}[t]{0.11\linewidth}
   \textcolor{green}{\texttt{{\refstepcounter{cntr}\label{#1}In \thecntr:}}}
  \end{minipage}
  \begin{minipage}[t]{0.89\linewidth}
  \begin{minted}[
    breaklines=true,style=monokai]{#1}}
 {\end{minted}
 \end{minipage}}

\newenvironment{pythonout}[1]
{\VerbatimEnvironment
  \addtocounter{cntr}{-1}
  \begin{minipage}[t]{0.11\linewidth}
   \textcolor{red}{\texttt{{\refstepcounter{cntr}\label{#1}Out\thecntr:}}}
  \end{minipage}
  \begin{minipage}[t]{0.89\linewidth}
  \begin{minted}[
    breaklines=true,style=monokai]{#1}}
 {\end{minted}
 \end{minipage}}

\newenvironment{pythonerr}[1]
{\VerbatimEnvironment
  \begin{minted}[
    breaklines=true,style=monokai]{#1}}
 {\end{minted}}

\begin{document}

% This disables the \pause command, handy in the editing phase.
%\renewcommand{\pause}{}

% Make the title slide.
\makeTitleSlide{\includegraphics[height=3.5cm]{../../images/Python.pdf}}

% First page of the presentation.
\section{Introduction}
\makeTableOfContents


\section{Sequence types}

\subsection{Lists}
\begin{pframe}
Mutable sequences of values.
 \begin{ipython}
  \begin{pythonin}{python}
l = [2, 5, 2, 3, 7]
  \end{pythonin}
  \\
  \begin{pythonin}{python}
type(l)
  \end{pythonin}
  \begin{pythonout}{python}
list
  \end{pythonout}
 \end{ipython}
 \medskip

 Lists can be heterogeneous, but we typically don't use that.
 \begin{ipython}
  \begin{pythonin}{python}
a = 'spezi'
  \end{pythonin}
  \\
  \begin{pythonin}{python}
[3, 'abc', 1.3e20, [a, a, 2]]
  \end{pythonin}
  \begin{pythonout}{python}
[3, 'abc', 1.3e+20, ['spezi', 'spezi', 2]]
  \end{pythonout}
 \end{ipython}
\end{pframe}


\subsection{Tuples}
\begin{pframe}
Mutable sequences of values.
 \begin{ipython}
  \begin{pythonin}{python}
t = 'white', 77, 1.5
  \end{pythonin}
  \\
  \begin{pythonin}{python}
type(t)
  \end{pythonin}
  \begin{pythonout}{python}
typle
  \end{pythonout}
 \\

 \begin{pythonin}{python}
color, width, scale = t
  \end{pythonin}
  \\
  \begin{pythonin}{python}
width
  \end{pythonin}
  \begin{pythonout}{python}
77
  \end{pythonout}
 \end{ipython}
\end{pframe}


\subsection{Strings}
\begin{pframe}
Immutable sequences of characters.
 \begin{ipython}
  \begin{pythonin}{python}
'a string can be written in single quotes'
  \end{pythonin}
  \\
  \begin{pythonout}{python}
'a string can be written in single quotes'
  \end{pythonout}
 \end{ipython}

Strings can also be written with double quotes, or over multiple lines with
triple-quotes.
 \begin{ipython}
  \begin{pythonin}{python}
"this makes it easier to use the ' character"
  \end{pythonin}
  \\
  \begin{pythonout}{python}
"this makes it easier to use the ' character"
  \end{pythonout}
  \\

  \begin{pythonin}{python}
"""A multiline string.
You see? I continued after a blank line."""
  \end{pythonin}
  \begin{pythonout}{python}
'A multiline string.\n\nYou see? I continued after a blank line.'
  \end{pythonout}
 \end{ipython}
\end{pframe}

\begin{pframe}
 A common operation is formatting strings using argument substitutions.
 \begin{ipython}
  \begin{pythonin}{python}
'{} times {} equals {:.2f}'.format('pi', 2, 6.283185307179586)
  \end{pythonin}
  \\
  \begin{pythonout}{python}
'pi times 2 equals 6.28'
  \end{pythonout}
 \end{ipython}

 Accessing arguments by position or name is more readable.
 \begin{ipython}
  \begin{pythonin}{python}
'{1} times {0} equals {2:.2f}'.format('pi', 2, 6.283185307179586)
  \end{pythonin}
  \\
  \begin{pythonout}{python}
'2 times pi equals 6.28'
  \end{pythonout}
  \\

  \begin{pythonin}{python}
'{number} times {amount} equals {result:.2f}'.format(number='pi', amount=2, result=6.283185307179586)
  \end{pythonin}
  \\
  \begin{pythonout}{python}
'pi times 2 equals 6.28'
  \end{pythonout}
 \end{ipython}
\end{pframe}


\section{Common sequence operations}

\begin{pframe}
 All sequence types support: concatenation, membership/substring tests,
 indexing, and slicing.
 \\
 \begin{ipython}
  \begin{pythonin}{python}
[1, 2, 3] + [4, 5, 6]
  \end{pythonin}
  \\
  \begin{pythonout}{python}
[1, 2, 3, 4, 5, 6]
  \end{pythonout}
  \\

  \begin{pythonin}{python}
'bier' in 'we drinken bier vanaf half 5'
  \end{pythonin}
  \\
  \begin{pythonout}{python}
True
  \end{pythonout}
  \\

  \begin{pythonin}{python}
'abcdefghijkl'[5]
  \end{pythonin}
  \\
  \begin{pythonout}{python}
'f'
  \end{pythonout}
 \end{ipython}
\end{pframe}


\subsection{Slicing}
\begin{pframe}
 Slice \mintinline{python}{s} from \mintinline{python}{i} to
 \mintinline{python}{j} with \mintinline{python}{s[i:j]}.
 \begin{ipython}
  \begin{pythonin}{python}
'abcdefghijkl'[4:8]
  \end{pythonin}
  \\
  \begin{pythonout}{python}
'efgh'
  \end{pythonout}
  \\

  \begin{pythonin}{python}
'abcdefghijkl'[:3]
  \end{pythonin}
  \\
  \begin{pythonout}{python}
'abc'
  \end{pythonout}
 \end{ipython}

 We can also define the step \mintinline{python}{k} with
 \mintinline{python}{s[i:j:k]}.
 \begin{ipython}
  \begin{pythonin}{python}
'abcdefghijkl'[7:3:-1]
  \end{pythonin}
  \\
  \begin{pythonout}{python}
'hgfe'
  \end{pythonout}
 \end{ipython}
\end{pframe}


\subsection{Several helpful builtins}
\begin{pframe}
 \begin{ipython}
  \begin{pythonin}{python}
len('attacgataggcatccgt')
  \end{pythonin}
  \\
  \begin{pythonout}{python}
18
  \end{pythonout}
  \\

  \begin{pythonin}{python}
max([17, 86, 34, 51])
  \end{pythonin}
  \\
  \begin{pythonout}{python}
86
  \end{pythonout}
  \\

  \begin{pythonin}{python}
('atg', 22, True, 'atg').count('atg')
  \end{pythonin}
  \\
  \begin{pythonout}{python}
2
  \end{pythonout}
 \end{ipython}
\end{pframe}


\subsection{More with lists}
\begin{pframe}
 We can replace, add, remove, reverse and sort items in-place.
 \begin{ipython}
  \begin{pythonin}{python}
l = [1, 2, 3, 4]
  \end{pythonin}
  \\
  \begin{pythonin}{python}
l[3] = 7
  \end{pythonin}
  \\
  \begin{pythonin}{python}
l.append(1)
  \end{pythonin}
  \\
  \begin{pythonin}{python}
l[1:3] = [3, 2]
  \end{pythonin}
  \\
  \begin{pythonin}{python}
l.sort()
  \end{pythonin}
  \\
  \begin{pythonin}{python}
l.reverse()
  \end{pythonin}
  \\
  \begin{pythonin}{python}
l
  \end{pythonin}
  \\
  \begin{pythonout}{python}
[7, 3, 2, 1, 1]
  \end{pythonout}
 \end{ipython}
\end{pframe}



\section{Dictionaries}

\subsection{Dictionaries map hashable values to arbitrary objects}
\begin{pframe}
 \begin{itemize}
  \item All built-in immutable objects are hashable.
  \item No built-in mutable objects are hashable.
 \end{itemize}
 \begin{ipython}
  \begin{pythonin}{python}
d = {'a': 27, 'b': 18, 'c': 12}
  \end{pythonin}
  \\
  \begin{pythonin}{python}
type(d)
  \end{pythonin}
  \\
  \begin{pythonout}{python}
dict
  \end{pythonout}
  \\

  \begin{pythonin}{python}
d['e'] = 17
  \end{pythonin}
  \\
  \begin{pythonin}{python}
'e' in d
  \end{pythonin}
  \\
  \begin{pythonout}{python}
True
  \end{pythonout}
  \\

  \begin{pythonin}{python}
d.update({'a': 18, 'f': 2})
  \end{pythonin}
  \\
  \begin{pythonin}{python}
d
  \end{pythonin}
  \\
  \begin{pythonout}{python}
{'a': 18, 'b': 18, 'c': 12, 'e': 17, 'f': 2}
  \end{pythonout}
 \end{ipython}
\end{pframe}


\subsection{Accessing dictionary content}
\begin{pframe}
 \begin{ipython}
  \begin{pythonin}{python}
d['b']
  \end{pythonin}
  \\
  \begin{pythonout}{python}
18
  \end{pythonout}
  \\

  \begin{pythonin}{python}
d.keys()
  \end{pythonin}
  \\
  \begin{pythonout}{python}
dict_keys(['e', 'c', 'a', 'b', 'f'])
  \end{pythonout}
  \\

  \begin{pythonin}{python}
list(d.keys())
  \end{pythonin}
  \\
  \begin{pythonout}{python}
['a', 'c', 'b', 'e', 'f']
  \end{pythonout}
  \\

  \begin{pythonin}{python}
list(d.values())
  \end{pythonin}
  \\
  \begin{pythonout}{python}
[18, 12, 18, 17, 2]
  \end{pythonout}
  \\

  \begin{pythonin}{python}
list(d.items())
  \end{pythonin}
  \\
  \begin{pythonout}{python}
[('a', 18), ('c', 12), ('b', 18), ('e', 17), ('f', 2)]
  \end{pythonout}
 \end{ipython}
\end{pframe}



\section{Sets}

\subsection{Mutable unordered collections of hashable values without duplication}
\begin{pframe}
 \begin{ipython}
  \begin{pythonin}{python}
x = {12, 28, 21, 17}
  \end{pythonin}
  \\
  \begin{pythonin}{python}
type(x)
  \end{pythonin}
  \\
  \begin{pythonout}{python}
set
  \end{pythonout}
  \\

  \begin{pythonin}{python}
x.add(12)
  \end{pythonin}
  \\
  \begin{pythonin}{python}
a
  \end{pythonin}
  \\
  \begin{pythonout}{python}
{12, 17, 21, 28}
  \end{pythonout}
  \\

  \begin{pythonin}{python}
x.discard(21)
  \end{pythonin}
  \\
  \begin{pythonin}{python}
x
  \end{pythonin}
  \\
  \begin{pythonout}{python}
{12, 17, 28}
  \end{pythonout}
 \end{ipython}
\end{pframe}


\subsection{Operations}
\begin{pframe}
 We can test for membership and apply many common set operations\\
 such as union and intersect.
 \medskip

 \begin{ipython}
  \begin{pythonin}{python}
17 in {12, 28, 21, 17}
  \end{pythonin}
  \\
  \begin{pythonout}{python}
True
  \end{pythonout}
  \\

  \begin{pythonin}{python}
{12, 28, 21, 17} | {12, 18, 11}
  \end{pythonin}
  \\
  \begin{pythonout}{python}
{11, 12, 17, 18, 21, 28}
  \end{pythonout}
  \\

  \begin{pythonin}{python}
{12, 28, 21, 17} & {12, 18, 11}
  \end{pythonin}
  \\
  \begin{pythonout}{python}
{12}
  \end{pythonout}
 \end{ipython}
\end{pframe}



\section{Booleans}

\begin{pframe}
  The two boolean values are written \mintinline{python}{False} and
  \mintinline{python}{True}.
 \begin{ipython}
  \begin{pythonin}{python}
True or False
  \end{pythonin}
  \\
  \begin{pythonout}{python}
True
  \end{pythonout}
  \\

  \begin{pythonin}{python}
True and False
  \end{pythonin}
  \\
  \begin{pythonout}{python}
False
  \end{pythonout}
  \\

  \begin{pythonin}{python}
not False
  \end{pythonin}
  \\
  \begin{pythonout}{python}
True
  \end{pythonout}
 \end{ipython}
\end{pframe}


\subsection{Comparisons}
\begin{pframe}
  Comparisons can be done on all objects and return a boolean value.
  \mintinline{python}{True}.
 \begin{ipython}
  \begin{pythonin}{python}
22 * 3 > 66
  \end{pythonin}
  \\
  \begin{pythonout}{python}
False
  \end{pythonout}
 \end{ipython}

 We have two equivalence relations: value equality (\mintinline{python}{==}) and
 object identity (\mintinline{python}{is}).
 \begin{ipython}
  \begin{pythonin}{python}
a, b = [1, 2, 3], [1, 2, 3]
  \end{pythonin}
  \\
  \begin{pythonin}{python}
a == b
  \end{pythonin}
  \\
  \begin{pythonout}{python}
True
  \end{pythonout}
  \\

  \begin{pythonin}{python}
a is b
  \end{pythonin}
  \\
  \begin{pythonout}{python}
False
  \end{pythonout}
 \end{ipython}


\end{pframe}


\section{Hands on!}
\begin{pframe}
 \vspace{-0.5cm}
 \begin{enumerate}
  \item Make a list with 10 integer elements.
  \begin{enumerate}[a]
  \item What is the sum of all the items in the list.
  \item Make a new list from the above one that does not include the 0th, 4th,
  and 5th elements.
  \item Sum only the elements from the first list which are between the 2nd and
  the 6th elements.
  \item Make a new list that includes only the elements that are greater than
  $10$ from the first list.
  \end{enumerate}
  \item Food:
  \begin{enumerate}[a.]
   \item Create a dictionary for food products called "prices" and put some
   values in it, e.g., "apples": 2, "oranges": 1.5, "pears": 3, ...
   \item Create a corresponding dictionary called "stocks" and put the stock
   values in it, e.g., "apples": 0, "oranges": 1, "pears": 10, ...
   \item Print stock and price information for each food item.
   \item Determine and print how much money you would make if you sold all of
   your food products.
  \end{enumerate}
 \end{enumerate}
\end{pframe}


% Make the acknowledgements slide.
\makeAcknowledgementsSlide{
  \begin{tabular}{ll}
    Martijn Vermaat\\
    Jeroen Laros\\
    Jonathan Vis
  \end{tabular}
}

\end{document}
